
This work establishes two findings regarding benchmark fingerprints in fine-tuned vision models:

\textbf{Fingerprints are detectable with high accuracy.} A linear classifier identifies training benchmark origin from ResNet-50 penultimate layer features with 99.51\% accuracy (Cohen's d = 698.08). This demonstrates that fine-tuning encodes benchmark-specific information that is linearly separable in representation space.

\textbf{The BFS metric does not predict generalization gap under tested conditions.} The correlation between Benchmark Fingerprint Score and cross-dataset performance gap is r = 0.022 (p = 0.967). The proposed mechanism linking fingerprint strength to generalization degradation is not supported by these data, though ceiling effects in BFS (all models > 0.999) prevent conclusive interpretation.

The observed paradox---high detectability but no predictive power---points to questions about what fine-tuning changes in representations and how those changes relate to deployment performance. Future work should explore calibrated fingerprint metrics that provide variance across models, layer-wise analysis to identify where fingerprints emerge, and intervention studies (e.g., representation regularization) to establish causal relationships between fingerprint strength and generalization.

The 99.51\% classification accuracy demonstrates that models encode their training benchmarks. What remains to be established is whether and how that encoding relates to deployment performance.

